\documentclass[12pt]{article}
\usepackage[utf8]{inputenc}
\usepackage[T1]{fontenc}
\usepackage{hyperref}
\title{Vinza Language Resources for K-12 Curriculum Development}
\author{Compiled by Claude Sonnet 5 (claude-sonnet-5) for Andr\'as Kornai}
\date{}
\begin{document}
\maketitle

\section{Introduction}
Vinza (ISO 639-3 \texttt{vin}) was listed as a language of Tanzania near Uvinza. Ethnologue states it could find no evidence the language exists and is pursuing removal of the code from ISO 639-3; I did not find the formal retirement record. Older listings (Wikipedia, ASJP, Global Recordings) still carry it, with a 1987 figure of about 10,000 speakers. No curriculum resources can responsibly be listed.

\section{Primary Educational Resources}
None found.

\section{Linguistic and Reference Materials}
\begin{itemize}
\item Ethnologue entry for vin: \url{https://www.ethnologue.com/language/vin}
\item Vinza language (Wikipedia): \url{https://en.wikipedia.org/wiki/Vinza_language}
\item ASJP entry: \url{https://asjp.clld.org/languages/VINZA}
\item Endangered Languages Project entry: \url{https://endangeredlanguages.com/lang/5073}
\end{itemize}

\section{Religious and Community Resources}
None found. Joshua Project lists the code: \url{https://joshuaproject.net/languages/vin}

\section{Digital and Audio Resources}
Global Recordings Network lists the code (contents not verified): \url{https://globalrecordings.net/en/language/vin}

\section{Curriculum Development Context}
Since the language's existence is in doubt, a dedicated curriculum is not supportable. If a real community exists under this name, it should first be identified, and the neighbouring languages of the Uvinza area (and Swahili) would be the relevant educational languages.
\end{document}